% !TEX TS-program = xelatex
% !TEX encoding = UTF-8 Unicode
% -*- coding: UTF-8; -*-
% vim: set fenc=utf-8

%%%%%%%%%%%%%%%%%%%%%%%%%%%%%%%%%%%%%%%%%%%%%%%%%%%%%%%%%%%%%%%%%
%% SIMPLE-THESIS-DISSERTATION
%% <https://github.com/zachscrivena/simple-thesis-dissertation>
%% This is free and unencumbered software released into the
%% public domain; see <http://unlicense.org> for details.
%%%%%%%%%%%%%%%%%%%%%%%%%%%%%%%%%%%%%%%%%%%%%%%%%%%%%%%%%%%%%%%%%

% See "README.md" for instructions on compiling this document.
\documentclass[letterpaper,nonstop,draftmode]{simplethesisdissertation}
\usepackage{tikz}
\usepackage{tensor}
\usepackage{braket}
% Class options:
% a4paper, letterpaper, nonstop, draftmode.

%%%%%%%%%%%%%%%%%%%%%%%%%%%%%%%%%%%%%%%%%%%%%%%%%%%%%%%%%%%%%%%%%
%% PREAMBLE.
%%%%%%%%%%%%%%%%%%%%%%%%%%%%%%%%%%%%%%%%%%%%%%%%%%%%%%%%%%%%%%%%%

% Document properties.
\def\DocumentTitle{Insert Thesis Title Here}
\def\AuthorName{Insert Author Name Here}

% PDF settings and properties.
\hypersetup{
pdftitle={\DocumentTitle},
pdfauthor={\AuthorName},
pdfsubject={Ph.D. Thesis, University Institute of College, 2016},
pdfcreator={XeLaTeX},
pdfproducer={},
pdfkeywords={},
unicode=true,
bookmarks=true,
bookmarksopen=true,
pdfstartview=FitH,
pdfpagelayout=OneColumn,
pdfpagemode=UseOutlines,
hidelinks,
breaklinks,
bookmarksnumbered}

% Macros:
\DeclareMathOperator*{\argmax}{arg\,max}
\DeclareMathOperator*{\argmin}{arg\,min}
\renewcommand{\binom}[2]{\left(\genfrac{}{}{0pt}{}{#1}{#2}\right)}
\newcommand{\ceil}[1]{{\left\lceil{#1}\right\rceil}}
%\newcommand{\ffrac}[2]{{\nicefrac{#1}{#2}}}
%\newcommand{\fffrac}[2]{{\left.{#1}\middle/{#2}\right.}}
\newcommand{\floor}[1]{{\left\lfloor{#1}\right\rfloor}}
\DeclareMathOperator{\lcm}{lcm}
\newcommand{\ZZ}{{\mathbb{Z}}}

%%%%%%%%%%%%%%%%%%%%%%%%%%%%%%%%%%%%%%%%%%%%%%%%%%%%%%%%%%%%%%%%%
%% ACTUAL DOCUMENT.
%%%%%%%%%%%%%%%%%%%%%%%%%%%%%%%%%%%%%%%%%%%%%%%%%%%%%%%%%%%%%%%%%

\begin{document}

% Use Roman numerals (i, ii, iii, etc.) for page numbers in the front matter.
\pagenumbering{roman}

%%%%%%%%%%%%%%%%%%%%%%%%%%%%%%%%%%%%%%%%%%%%%%%%%%%%%%%%%%%%%%%%%
%% TITLE PAGE.
%%%%%%%%%%%%%%%%%%%%%%%%%%%%%%%%%%%%%%%%%%%%%%%%%%%%%%%%%%%%%%%%%

% No headers or footers on the title page.
\thispagestyle{empty}

\begingroup
\centering
\setstretch{1.0}
~
\\[1em]
\sffamily\bfseries\fontsize{26}{31.2}\selectfont
\DocumentTitle
\\[0.6in]
\sffamily\bfseries\Large
\AuthorName
\\[0.6in]
\normalfont\normalsize
An undergraduate thesis advised by Dr. David Craig \\~\\
submitted to the Department of Physics, Oregon State University \\~\\
in partial fulfillment of the requirements for the degree BSc in Physics \\[0.6in]
Submitted on May 8, 2020
\vfill
\includegraphics[height=1.0in]
{Figure-SchoolLogo}
\par
\endgroup

\clearpage

%%%%%%%%%%%%%%%%%%%%%%%%%%%%%%%%%%%%%%%%%%%%%%%%%%%%%%%%%%%%%%%%%
%% COPYRIGHT PAGE.
%%%%%%%%%%%%%%%%%%%%%%%%%%%%%%%%%%%%%%%%%%%%%%%%%%%%%%%%%%%%%%%%%

% \pagestyle{plain}
% \setcounter{page}{2}
%
% \begingroup
% \centering
% \setstretch{1.0}
% \null
% \vfill
% {\sffamily\textcopyright}~2016
% \\[0.5em]
% \AuthorName
% \\[0.5em]
% All Rights Reserved
% \par
% \endgroup
%
% \clearpage

% %%%%%%%%%%%%%%%%%%%%%%%%%%%%%%%%%%%%%%%%%%%%%%%%%%%%%%%%%%%%%%%%%
% %% DEDICATION PAGE.
% %%%%%%%%%%%%%%%%%%%%%%%%%%%%%%%%%%%%%%%%%%%%%%%%%%%%%%%%%%%%%%%%%
%
% \begingroup
% \centering
% \setstretch{1.0}
% ~
% \\[1in]
% \textit{Insert dedication here}
% \par
% \endgroup
%
% \clearpage

%%%%%%%%%%%%%%%%%%%%%%%%%%%%%%%%%%%%%%%%%%%%%%%%%%%%%%%%%%%%%%%%%
%% ACKNOWLEDGMENTS.
%%%%%%%%%%%%%%%%%%%%%%%%%%%%%%%%%%%%%%%%%%%%%%%%%%%%%%%%%%%%%%%%%

\chapter*{Acknowledgments}
\addcontentsline{toc}{chapter}{Acknowledgments}

I would like to thank Dr. David Craig, for guiding me through an engaging line of research, as well as Dr. David McIntyre, Dr. Elizabeth Gire, Dr. Corinne Manogue and Dr. Janet Tate for developing the quantum curriculum from which this thesis is rooted.

I would also like to thank all of those who gave me the time, space, and support I needed while writing this. This includes (but is certainly not limited to) my partner Brooke, my parents Joy and Kevin, my housemate Cheyanne, the staff of Interzone, and my friends Saskia, Rachel and Justin.

\clearpage

%%%%%%%%%%%%%%%%%%%%%%%%%%%%%%%%%%%%%%%%%%%%%%%%%%%%%%%%%%%%%%%%%
%% ABSTRACT.
%%%%%%%%%%%%%%%%%%%%%%%%%%%%%%%%%%%%%%%%%%%%%%%%%%%%%%%%%%%%%%%%%

\chapter*{Abstract}
\addcontentsline{toc}{chapter}{Abstract}

Standard quantum mechanics makes foundational assumptions to describe the measurement process. Upon interaction with a ``classical measurement apparatus'', a quantum system is subjected to postulated ``state collapse'' dynamics. We show that framing measurement around state collapse and ill-defined classical observers leads to interpretational issues, and artificially limits the scope of quantum theory. This motivates describing measurement as a unitary process instead. In the context of the Stern-Gerlach experiment, the measurement of an electron's spin angular momentum is explained as the entanglement of its spin and position degrees of freedom. Furthermore, the electron-environment interaction is also detailed as part of the measurement process. The environment plays the role of a record keeper, establishing the ``facts of the universe'' to make the measurement's occurrence objective. This enables the assignment of probabilities to outcomes of consecutive Stern-Gerlach experiments, which is done using the consistent histories approach to quantum mechanics. We exemplify how consistent histories can be used to give a unitary measurement model predictive power by reproducing the predictions of the standard approach. In doing so, we provide a ``spins-first'' introduction to the consistent histories approach that is targeted towards anybody looking to learn more about quantum foundations. Along the way, we coin a mathematically convenient operator that simplifies the calculation of probabilities. We conclude by discussing how consistent histories ignores the ``problem of outcomes'' in quantum measurement theory, and the connection between this problem and the many worlds interpretation of quantum mechanics.

% A novel approach to limiting extraneous predictions implements the principle of complementarity into the measurement process directly.

\clearpage

%%%%%%%%%%%%%%%%%%%%%%%%%%%%%%%%%%%%%%%%%%%%%%%%%%%%%%%%%%%%%%%%%
%% TABLE OF CONTENTS (TOC), LISTS OF FIGURES, TABLES, ETC.
%%%%%%%%%%%%%%%%%%%%%%%%%%%%%%%%%%%%%%%%%%%%%%%%%%%%%%%%%%%%%%%%%

\tableofcontents

\listoffigures

\clearpage

% Use Arabic numerals (1, 2, 3, etc.) for subsequent page numbers.
\pagenumbering{arabic}


\part{Background}
\chapter{Stern-Gerlach Experiments}
\chapter{Conistent Histories Quantum Mechanics}


\usetikzlibrary{shapes.geometric}
\usetikzlibrary{positioning}

\newcommand{\apparatus}[4]{\node[square node] (#1) at (#2,#3){#4};
                           \node[port] (#1+) at (#2 + 0.375, #3 + 0.5){+};
                           \node[port] (#1-) at (#2 + 0.375, #3 - 0.5){-};}

\part{Stern-Gerlach Experiments}

\chapter{Postulates}

To discuss Stern-Gerlach experiments, we first consider the mathematical objects used to model physical systems and variables. By comparing the objects used in classical and quantum mechanics, we make sense of the first three postulates of quantum mechanics. Then, we compare standard and consistent descriptions of spin measurement, and their relation to the fourth and fifth postulates.

\section{State Spaces and Physical Variables}
Consider the spin (intrinsic angular momentum) of an electron. Our system is the electron, and the classical spin state is simply a vector $S = (S_x, S_y, S_z)$ in $\mathbb{R}^3$, where each component $S_w$ is a physical variable that represents the magnitude of spin in the $w$ direction. Consequently, the components of spin state are mutually independent, and they can take on any real value. In other words, the sample spaces for each component are compatible, continuous and infinitely large. Because of the compatibility of spin component sample spaces, succesive observations of $S_x$, $S_y$, and $S_z$ uniquely determine the spin state.

Our physical variable of interest is the magnitude of spin in the $z$ direction; clasically, this variable is modeled by a function of spin state: $S_z(S) = \braket{z|S}$. This function relies on the inner product of the state space $\mathbb{R}^3$ to determine the magnitude of the spin vector in the $z$ direction only. In general, measurement of spin in three orthogonal directions uniquely determines the spin component in any direction $w$.

Experimental evidence raises issues with this classical description of electron spin states and $S_z$. First, only two values have ever been recorded, named \textit{spin-up} and \textit{spin-down}. These values are of equal magnitude $\frac{\hbar}{2}$ with positive and negative orientations, respectively. Second, the results of succesive measurements of a spin system impliy that spin state is determined by spin along one axis only. Consequently, spin along more than one axis cannot be simultaneously determined; this is discussed further in Section \sref{Section:Measurement}. To resolve these issues, we must change the mathematical objects used to represent system states and physical variables. Specifically, the sample space must restrict observed $S_w$ values to spin up and spin down, and the inability of determining spin in more than one direction should be reflected through incompatible sample spaces.

The quantum mechanical system state is described by a normalized vector in a linear state space. In fact, this is the first postulate of quantum mechanics:
\begin{Thm:Postulate}{1}
    The state of a physical system is defined by specifying an abstract vector $\ket{\psi}$ in a Hilbert state space $\mathcal{H}$.
\end{Thm:Postulate}

For spin-$\frac{1}{2}$ systems, the state space consists of all linear combinations of spin-up and spin-down:
\begin{equation}
    \mathcal{H} = \{\alpha\ket{+} + \beta\ket{-}\}
\end{equation}
where $\alpha, \beta \in \mathbb{C}$

For ease of calculation, we only consider normalized states, as scalar multiples of any normalized state bear the same physical meaning. We require $\alpha^2 + \beta^2 = 1$.

Defining system states in this way has dramatic consequences. A state includes all information that can be known about a system. A classical spin state includes spin magnitude in $x$, $y$, and $z$. The quantum state is only capable of determining spin magintude in one general direction $w$; spin in other directions not parallel or antiparallel to $w$ is undefined.

The second posulate of quantum mechanics states that physical variables are described by linear operators:
\begin{Thm:Postulate}{2}
    Every physical variable $\mathcal{A}$ is described by an operator $A$ acting in $\mathcal{H}$.
\end{Thm:Postulate}

Justifying the second postulate is easiest when also considering the third postulate:
\begin{Thm:Postulate}{3}
    The only possible result of the measurement of a physical variable $\mathcal{A}$ is one of the eigenvalues of the corresponding operator $A$.
\end{Thm:Postulate}

The operator correlates elements of a finite sample space (eigenvalues) with particular system states (eigenstates). Consequently, a state can only be interpreted as having a definite variable value if it is an eigenstate of the variable's operator. To illustrate this, consider the operator representing $S_z$. Written in its own basis,
\begin{align*}
    S_z \doteq \frac{\hbar}{2}\begin{bmatrix} 1 & 0 \\ 0 & -1 \end{bmatrix}
\end{align*}
This operator correlates $z$ spin-up $\left(S_z = \frac{\hbar}{2}\right)$ with
\begin{align*}
    \ket{\psi} = \ket{+}_z \doteq \begin{bmatrix} 1 \\ 0 \end{bmatrix}
\end{align*}
and $z$ spin-down $\left(S_z = \frac{-\hbar}{2}\right)$ with
\begin{align*}
    \ket{\psi} = \ket{-}_z \doteq \begin{bmatrix} 0 \\ 1 \end{bmatrix}
\end{align*}
Similarly, the operator representing $S_y$ written in the $S_z$ basis is
\begin{align*}
        S_y \doteq \frac{\hbar}{2}\begin{bmatrix} 0 & -i \\ i & 0 \end{bmatrix}
\end{align*}
This operator correlates $y$ spin-up $\left(S_y = \frac{\hbar}{2}\right)$ with
\begin{align*}
    \ket{\psi} = \ket{+}_y \doteq \frac{1}{\sqrt{2}}\begin{bmatrix} 1 \\ i \end{bmatrix}
\end{align*}
and $y$ spin-down $\left(S_y = \frac{-\hbar}{2}\right)$ with
\begin{align*}
    \ket{\psi} = \ket{-}_y \doteq \frac{1}{\sqrt{2}}\begin{bmatrix} 1 \\ -i \end{bmatrix}
\end{align*}

Operators for $S_z$ and $S_y$ share no common eigenstates, so no state can posses definite values for both variables. In general, operators for any spin component $S_j$, where $j$ is some direction not parallel or antiparallel with $z$, share no common eigenstates with $S_z$; in other words, $S_z$ and $S_j$ have incompatible sample spaces. By representing physical variables with an operator rather than a function, sample spaces become quantized and can be incompatible with each other. These features are necessary for predicting experimental results.

The first three postulates designate the mathematical objects used to model physical system states and variables. Consistent quantum theory does not modify these postulates; we will see that all of the strange features unique to quantum systems are a consequence of these postulates.

\section{Measurement}
\label{Section:Measurement}
\subsection{Standard Description}
In standard quantum mechanics, the fifth postualte (known as the projection postulate) describes how a system changes upon measurement. Defined as ``an interaction with a classical apparatus'', measurement instantaneously changes the measured state to the eigenstate corresponding to the measruement result.

\begin{Thm:Postulate}{5}
    If the measurement of the physical variable $\mathcal{A}$ on the system in the state $\ket{\psi}$ gives the result $a_n$, the state of the system immediately after the measurement is the normalized projection
    \begin{align*}
        \ket{\psi}^\prime = \frac{P_n\ket{\psi}}{\sqrt{\bra{\psi}P_n\ket{\psi}
        }}
    \end{align*}
    onto the subspace associated with $a_n$.
\end{Thm:Postulate}

${P^n}_z = \ket{n}_z\tensor[_z]{\bra{n}}{}$ is the projection operator for the state $\ket{n}_z$ corresponding to $n_z$. Consider a measurement result for the $z$ component of spin, $S_z$. We represent the result with $n_z$, which could be either spin-up or spin-down. The new state is the normalized projection of $\ket{\psi}$ onto $\ket{n}_z$. In other words, $\ket{\psi}$ instantaneously becomes $\ket{n}_z$ upon measurement. This process is known as \textit{state collapse} or \textit{wavefunction collapse}.

\begin{figure}
\centering\CaptionFontSize
\begin{tikzpicture}[shorten >=1pt,auto, thick,
     square node/.style={rectangle, minimum height=2cm, minimum width=1.50cm, text width = 1.25cm, draw, font=\sffamily\Large\bfseries},
     port/.style={rectangle, draw,  minimum height=1cm, minimum width=0.75cm, font=\sffamily\Large\bfseries},
     wf/.style={rectangle, minimum height=1cm}]
    \apparatus{1}{3}{0}{$Z$};
    \apparatus{2}{6}{1.5}{$X$};

    \node[wf] (w0) at (0,0) {$\ket{\psi}$};
    \node[wf] (w1) at (8, 2.0) {$\ket{+}_x$};
    \node[wf] (w2) at (8, 1.0) {$\ket{-}_x$};

    \draw[line width=0.5mm] (w0) -- (1);

    \draw[line width=0.5mm] (1+) -- (2) node [near end] {$\ket{+}_z$};

    \draw[line width=0.5mm] (2+) -- (w1);
    \draw[line width=0.5mm] (2-) -- (w2);
\end{tikzpicture}
\caption[Insert an abbreviated caption here to show in the List of Figures]
{Demonstrating renormalizing upon measurment in standard quantum mechanics}
\label{Figure:Measurement:Renormalizing}
\end{figure}

As an example, consider the system shown in \fref{Figure:Measurement:Renormalizing}. The first apparatus serves as a state preparation device, since we are only interested in particles exiting the spin-up output. Using the projection postulate, the state after the first measurement is
\begin{align*}
    \ket{\psi_{top}} &= \frac{{P^z}_+\ket{\psi}}{\sqrt{\bra{\psi}{P^z}_+\ket{\psi}}} = \ket{+}_z
\end{align*}
Similarly, the possible output states from the second apparatus are $\tensor[_z]{\bra{+}}{}$
\begin{align*}
    \ket{\psi_{top}} &= \frac{{P^x}_+\ket{+}_z}{\sqrt{\tensor[_z]{\bra{+}}{}{P^x}_+\ket{+}_z}} = \ket{+}_x \\ \\
    \ket{\psi_{bottom}} &= \frac{{P^x}_-\ket{+}_z}{\sqrt{\tensor[_z]{\bra{+}}{}{P^x}_-\ket{+}_z}} = \ket{-}_x
\end{align*}

\subsection{Consistent Description}
In consistent quantum theory, what is meant by ``measurement'' in the projection postulate is itself modeled as a physical process. Therefore, no notion of an undefined ``interaction with a classical apparatus'' is required to describe how states evolve when their classical properties are recorded. To describe measurement, each Stern-Gerlach apparatus has its own detector state space, containing orthonormal states representing each measurement result. For a Stern-Gerlach apparatus, a detection state is defined by a particle existing at some spatially separated output. In general, a detection state is some classical indicator; examples include an apparatus needle pointing up, or a particle colliding with a screen in a distinguishable region.

Let the state space of the $z$ apparatus be represented by
\begin{align*}
    {D}_z =  \{\ket{D_+}_z, \ket{D_-}_z\}
\end{align*}
where
\begin{align*}
    _z\braket{D_+|D_+}_z&=1 \\
    _z\braket{D_+|D_-}_z&=0
\end{align*}
The requirement of orthonormality reflects mutually exlcusive outcomes; the particle exits spin-up or spin-down. The state space for an apparatus aligned in any general direction ${D}_w$ is similarly defined.

The act of measurement is described by correlating detector states with $\ket{\psi}$ eigenstates. The system then evolves unitarily by
\begin{align}
    \nonumber V: \\
    & \ket{\psi} = \left(\sum_{n} {P^z}_n\ket{\psi}\right) \otimes \left(\sum_{m} P_m\ket{D}\right) \mapsto \sum_{n} {P^z}_n\ket{\psi} \otimes \ket{D_n}
\end{align}
where ${P^z}_n$ is the projection operator for the $n^{th} z$ spin eigenstate and $P_m$ is the projection operator for the $m^{th}$ detector state.

Before, $\ket{\psi}$ was a superposition of tensor products of any spin eigenstate and any detector state. Any detector state could be realized with any quantum state. Afterwards, $\ket{\psi}$ is a superposition of tensor products of a spin eigenstate and one specific detector state; a detector state is realized with one specific quantum state only.

Figure \fref{Figure:Measurement:DetectorStates} illustrates this description of measurement. As physical variables with incompatible sample spaces are recorded ($S_z$ and $S_x$), $\ket{\psi}$ loses information about its history; by posessing a definite $S_x$ value, nothing can be said about $S_z$. However, the detector states are not impacted by these measurements, and the \textit{event history} of the state is preserved. Our system now belongs to a \textit{history Hilbert space} defined by $\widetilde{\mathcal{H}} = \mathcal{H} \otimes D_z \otimes D_x$.

Each pure state in $\widetilde{\mathcal{H}}$ represents a \textit{history} (or sequence of events). Because of the orthonormality conditions imposed on detector states, these histories satisfy \textit{consistency conditions}; that is, they are disjoint and their magnitudes sum to 1. Together, these pure states are a \textit{consistent family of histories}, meaning that they constitue an exhaustive and disjoint set of event sequences.

TODO: introduce another example that illustrates +x, +z interfere, but their histories do not.

\begin{figure}
\centering\CaptionFontSize
\begin{tikzpicture}[shorten >=1pt,auto, thick,
     square node/.style={rectangle, minimum height=2cm, minimum width=1.50cm, text width = 1.25cm, draw, font=\sffamily\Large\bfseries},
     port/.style={rectangle, draw,  minimum height=1cm, minimum width=0.75cm, font=\sffamily\Large\bfseries},
     wf/.style={rectangle, minimum height=1cm, text width = 0.7cm}]
    \apparatus{1}{3}{0}{$Z$};
    \apparatus{2}{6}{1.5}{$X$};

    \node[wf] (w0) at (0,0) {$\ket{\psi}$};
    \node[wf] (w1) at (8, 2.0) {$\ket{+}_x \otimes \ket{D_z}_+  \otimes \ket{D_x}_+$};
    \node[wf] (w2) at (8, 1.0) {$\ket{-}_x \otimes \ket{D_z}_+ \otimes \ket{D_x}_-$};

    \draw[line width=0.5mm] (w0) -- (1);

    \draw[line width=0.5mm] (1+) -- (2) node [near end] {$\ket{+} \otimes \ket{D_z}_+$};

    \draw[line width=0.5mm] (2+) -- (w1);
    \draw[line width=0.5mm] (2-) -- (w2);
\end{tikzpicture}
\caption[Insert an abbreviated caption here to show in the List of Figures]
{Demonstrating description of measurment as an abstract physical process in consistent quantum mechanics}
\label{Figure:Measurement:DetectorStates}
\end{figure}

Measurment is now described by an abstract physical process. In this model, it is now feasible for state collapse to occur independent of physicists conducting clever experiments. For this sytem, we defined a detector state as a Stern Gerlach apparatus recording spin up or spin down. However, one could define detector states as any record of quantum systems possesing physical properties that interact with the classical world accordingly. Detector states describe a system's behavior classically without necessarily including any classical measurement apparatus.

\section{The Born Rule}
\subsection{Standard Description}
The fourth postualte (known as the probability postulate) assigns a probability distribution to the sample space of a physical variable.

\begin{Thm:Postulate}{4}
    The probability $\mathcal{P}(n)$ of obtaining a result $a_n$ corresponding to $\ket{a}_n$  is equal to
     \begin{align*}
        \mathcal{P}(n) = |\tensor[_n]{\braket{a|\psi}}{}|^2
    \end{align*}
\end{Thm:Postulate}

This postulate is also known as the \textit{Born Rule}, which is usually presented in the language of wavefunctions. The probability that a system is found at position $x$ is the magnitude of the wavefunction at that point:
\begin{align}
    \mathcal{P}_x = |\psi(x)|^2
\end{align}

The probabilities assigned to each state leaving the $S_x$ Stern-Gerlach device in Figure \fref{Figure:Measurement:DetectorStates} are
\begin{align*}
    \mathcal{P}_{+_x} &= |_x\braket{+|+}|^2 = \frac{1}{2} \\
    \mathcal{P}_{-_x} &= |_x\braket{-|+}|^2 = \frac{1}{2}
\end{align*}
\subsection{Consistent Description}
The spirit of the Born Rule is unchanged in consistent quantum theory. The difference is that we now compute the inner product of the time evolved state, $V\ket{\psi}$, and the detection state representing the result we are interested in, $\ket{D_n}$. So, we assign probabilities with the following:

The probability $\mathcal{P}(n)$ of observing a result $a_n$ at a detector $D$ is equal to
\begin{align*}
  \mathcal{P}(n) = |\braket{D_n|V|\psi}|^2
\end{align*}

By taking the innner product with a detector state, we sum the probability amplitudes of all the histories that include that detection state. Histories that include the other spin outcome are orthogonal, and contribute nothing to this sum.

\subsection{Density Matrices}
The Born Rule can also be expressed in terms of \textit{density matrices}. From (TODO: reference appendix), we know that an inner product can be written as the trace of the corresponding dyad. The Born Rule is the complex square of an inner product, so we should be able to assign probabilities to detection states by tracing over some corresponding tensor. We will manipulate the Born Rule to take this form, and then examine the resulting tensor.

Expanding the complex square,
\begin{align*}
  \mathcal{P}(n) &= \left(\bra{D_n}V\ket{\psi}\right) \cdot \left(\bra{D_n}V\ket{\psi}\right)^* \\
  \mathcal{P}(n) &= \left(\bra{D_n}V\ket{\psi}\right) \cdot \left(\bra{\psi}V^\dagger\ket{D_n}\right)
\end{align*}
Rewriting the second inner product as the trace of a dyad,
\begin{align*}
  \mathcal{P}(n) &= \left(\bra{D_n}V\ket{\psi}\right) \cdot Tr\left(\ket{D_n}\bra{\psi}V^\dagger\right)f
\end{align*}
The remaining inner product, like any, is a scalar. Since trace is a linear operator, we can scale any factor inside the operation by this inner product.
\begin{align*}
  \mathcal{P}(n) &= Tr\left(\ket{D_n} \cdot \left(\bra{D_n}V\ket{\psi}\right) \cdot \left(\bra{\psi}V^\dagger\right)\right)
\end{align*}
Now we can rewrite $\ket{D_n}\bra{D_n}$ as ${P^D}_n$, and simplify grouping:
\begin{align*}
  \mathcal{P}(n) &= Tr\left({P^D}_n \cdot \left(V\ket{\psi}\bra{\psi}V^\dagger\right)\right)
\end{align*}

TODO: discuss remaining object ${P^D}_n \cdot \left(V\ket{\psi}\bra{\psi}V^\dagger\right)$

\chapter{Complementarity}

We now discuss the the principle of complementarity in the context of standard and consistent quantum mechanics. Arguably the most fundamental feature of quantum mechanics, the principle of complementarity states that a quantum system has pairs of physical observables which cannot be measured simultaneously. Components of spin on orthogonal axes are complementary properties, so we examine measurements of succesive Stern-Gerlach experiments.

First, we compute the probabilities of each outcome using standard quantum mechanics. The first apparatus serves as a state preparation device with output $\ket{+}$. By the direction of the projection postulate, the state is renormalized upon each measurement. After measuring a property complementary to what is known (such as spin along $x$, knowing spin along $z$), any information known about the input state is lost; the input state instantaneously changes to the state corresponding to the observed quantity. Consequently, there is an equal probability of observing the final state as $\ket{+}$ or $\ket{-}$ at either final apparatus, even though the state was initially prepared as $\ket{+}$, since
\begin{align*}
    \mathcal{P}_n &= |\braket{+|+}_x|^2 \\
                  &= |\braket{-|+}_x|^2 \\
                  &= |\braket{+|-}_x|^2 \\
                  &= |\braket{-|-}_x|^2 \\
                  &= \frac{1}{4}
\end{align*}
This contradiction with classical intuition is a direct result of the projection postulate. The act of measurement causes the system to shed properties previously recorded. TODO: describe measurement problem.

TODO: Demonstrate complementarity and single framework rule

\begin{figure}
\centering\CaptionFontSize
\begin{tikzpicture}[shorten >=1pt,auto, thick,
     square node/.style={rectangle, minimum height=2cm, minimum width=1.50cm, text width = 1.25cm, draw, font=\sffamily\Large\bfseries},
     port/.style={rectangle, draw,  minimum height=1cm, minimum width=0.75cm, font=\sffamily\Large\bfseries},
     wf/.style={rectangle, minimum height=1cm}]
    \apparatus{1}{3}{0}{$Z$};
    \apparatus{2}{6}{1.25}{$X$};
    \apparatus{3}{9}{2.50}{$Z$};
    \apparatus{4}{9}{0}{$Z$};

    \node[wf] (w0) at (0,0) {$\ket{\psi}$};
    \node[wf] (w1) at (11, 3.0) {$\ket{+}$};
    \node[wf] (w2) at (11, 2.0) {$\ket{-}$};
    \node[wf] (w3) at (11, 0.5) {$\ket{+}$};
    \node[wf] (w4) at (11, -0.5) {$\ket{-}$};

    \draw[line width=0.5mm] (w0) -- (1);

    \draw[line width=0.5mm] (1+) -- (2) node [near end] {$\ket{+}$};

    \draw[line width=0.5mm] (2-) -- (4) node [midway, below] {$\ket{-}_x$};
    \draw[line width=0.5mm] (2+) -- (3) node [midway, above] {$\ket{+}_x$};

    \draw[line width=0.5mm] (3-) -- (w2);
    \draw[line width=0.5mm] (3+) -- (w1);
    \draw[line width=0.5mm] (4-) -- (w4);
    \draw[line width=0.5mm] (4+) -- (w3);
\end{tikzpicture}
\caption[Insert an abbreviated caption here to show in the List of Figures]
{Demonstrating complementary measurments in standard quantum mechanics}
\label{Figure:Intro:FigureExampleC}
\end{figure}

\begin{figure}
\centering\CaptionFontSize
\begin{tikzpicture}[shorten >=1pt,auto, thick,
     square node/.style={rectangle, minimum height=2cm, minimum width=1.50cm, text width = 1.25cm, draw, font=\sffamily\Large\bfseries},
     port/.style={rectangle, draw,  minimum height=1cm, minimum width=0.75cm, font=\sffamily\Large\bfseries},
     wf/.style={rectangle, minimum height=1cm}]
    \apparatus{1}{3}{0}{$Z$};
    \apparatus{2}{6}{1.25}{$X$};
    \apparatus{3}{9}{2.50}{$Z$};
    \apparatus{4}{9}{0}{$Z$};

    \node[wf] (w0) at (0,0) {$\ket{\psi}$};
    \node[wf] (w1) at (11, 3.0) {$\ket{+}$};
    \node[wf] (w2) at (11, 2.0) {$\ket{-}$};
    \node[wf] (w3) at (11, 0.5) {$\ket{+}$};
    \node[wf] (w4) at (11, -0.5) {$\ket{-}$};

    \draw[line width=0.5mm] (w0) -- (1);

    \draw[line width=0.5mm] (1+) -- (2) node [near end] {$\ket{+}$};

    \draw[line width=0.5mm] (2-) -- (4) node [midway, below] {$\ket{-}_x$};
    \draw[line width=0.5mm] (2+) -- (3) node [midway, above] {$\ket{+}_x$};

    \draw[line width=0.5mm] (3-) -- (w2);
    \draw[line width=0.5mm] (3+) -- (w1);
    \draw[line width=0.5mm] (4-) -- (w4);
    \draw[line width=0.5mm] (4+) -- (w3);
\end{tikzpicture}
\caption[Insert an abbreviated caption here to show in the List of Figures]
{Demonstrating complementary measurments in consistent quantum mechanics}
\label{Figure:Intro:FigureExampleD}
\end{figure}

\begin{figure}
\centering\CaptionFontSize
\begin{tikzpicture}[shorten >=1pt,auto, thick,
     square node/.style={rectangle, minimum height=2cm, minimum width=1.50cm, text width = 1cm, draw, font=\sffamily\Large\bfseries},
     port/.style={rectangle, draw,  minimum height=1cm, minimum width=0.75cm, font=\sffamily\Large\bfseries},
     wf/.style={rectangle, minimum height=1cm}]
    \apparatus{1}{3}{0}{Z};
    \apparatus{2}{6}{1}{X};
    \apparatus{3}{9}{1}{Z};

    \node[wf] (w0) at (0,0) {$\ket{\psi}$};
    \node[wf] (w1) at (12,2.5) {$\ket{\psi_1}$};
    \node[wf] (w2) at (12,1.5) {$\ket{\psi_2}$};
    \node[wf] (w3) at (12,0.5) {$\ket{\psi_3}$};
    \node[wf] (w4) at (12,-0.5) {$\ket{\psi_4}$};

    \draw[line width=0.5mm] (w0) -- (1);

    \draw[transform canvas={yshift=-0.6em}, line width=0.5mm, loosely dotted] (1+) -- (2);
    \draw[transform canvas={yshift=-0.2em}, line width=0.5mm, dotted] (1+) -- (2);
    \draw[transform canvas={yshift=0.2em}, line width=0.5mm, dashed] (1+) -- (2);
    \draw[transform canvas={yshift=0.6em}, line width=0.5mm] (1+) -- (2);

    \draw[transform canvas={yshift=-0.4em}, line width=0.5mm, loosely dotted] (2-) -- (3);
    \draw[transform canvas={yshift=0em}, line width=0.5mm, dotted] (2-) -- (3);
    \draw[transform canvas={yshift=0em}, line width=0.5mm, dashed] (2+) -- (3);
    \draw[transform canvas={yshift=0.4em}, line width=0.5mm] (2+) -- (3);

    \draw[line width=0.5mm, loosely dotted] (3-) -- (w4);
    \draw[line width=0.5mm, dashed] (3-) -- (w3);
    \draw[line width=0.5mm, dotted] (3+) -- (w2);
    \draw[line width=0.5mm] (3+) -- (w1);
\end{tikzpicture}
\caption[Insert an abbreviated caption here to show in the List of Figures]
{TODO: create section to discuss this example, and how it creates an inconsistent set of histories. Describe how the set can be made to be consistent}
\label{Figure:Intro:FigureExampleE}
\end{figure}


\end{document}
